\scanpage{23}
\mspage{11}

4. \uline{Relations avec le cube orienté}

Si $I$ est un ens. à $4$ éléments, $J\simeq\mathscr P_{2,2}(I)\simeq V_I^*$ l'ens à trois éléments associé, on considère l'application de degré $2$

\diagram{\begin{tikzcd}[ampersand replacement=\&,column sep=small,row sep=0pt]
(*)\&\operatorname{Circ}(I)\arrow[r]\&\unsure{J|\mathbb J}.\\
\&\substack{\text{ens des}\\\text{permutations}\\\text{circulaires}\\\text{sur }I}\&u\mapsto u^2
\end{tikzcd}}

[ \inserted{NB} On a aussi l'application de degré $2$

\diagram{\begin{tikzcd}[ampersand replacement=\&,column sep=small,row sep=0pt]
\mathscr P_2(I)\arrow[r]\&J\\
\substack{\simeq\text{ens. des}\\\text{``transpositions}\\\text{de }I\\\text{(ou ``réflexions'')}}\&
\substack{A\mapsto\{A,I\smallsetminus A\}\\\text{ou }\forall\,\tau\ \text{réflexion,}\\\exists!\,\tau'\ \text{réflexion telle}\\\text{que }\tau\tau'\in V^*}
\end{tikzcd}}

ce sont vraiment des \unsure{revêt\deleted{ement}s|revêt\deleted{ts}s} de nature distincte, les espaces homogènes \inserted{sous $G$} $\operatorname{Circ}(I)$ et $\mathscr P_2(I)$ n'étant pas $\widetilde{\mathrm m}$ isomorphes (le stabilisateur d'un élément de $\operatorname{Circ}$ est $\simeq\mathbb Z/4\mathbb Z$, celui d'un élément de $\mathscr P_2(I)$ est $\simeq\mathbb F_2\times\mathbb F_2$) -- c'est le premier \deleted{des} deux, \inserted{\deleted{\unsure{$\mathscr P$|$\mathfrak P$}} $\operatorname{Circ}(I)$,} qui peut s'interpréter comme \deleted{l'application} \inserted{l'ensemble} \deleted{\unsure{de l'appl.}} \inserted{des} faces d'un cube\ldots\ ].

L'application de degré $2$ $(*)$ définit un cube \inserted{comb.}, dont l'ens. des faces est $\operatorname{Circ}(I)$, \inserted{avec l'antipodisme $u\mapsto u^{-1}$,} l'ens des \deleted{sommets} \deleted{\unsure{arêtes}} arêtes \inserted{(\unsure{resp$^{\mathrm t}$|compr$^{\mathrm t}$|compl.} des sommets)} est l'ens des $12$ sections partielles \deleted{\unsure{dans}} sur des parties de card. $2$ de $J$ (resp. des \deleted{\unsure{$8$|$\mathcal P$} \unsure{les}} $8$ sections de $\operatorname{Circ}(I)$ sur $J$), l'incidence entre facettes \deleted{\unsure{étant}} de dim $2,1,0$ étant le prolongement des \inserted{sections partielles.} \deleted{\unsure{À tout}} Il y aurait

% slanted
\margin{Bien sûr, $G$ opère sur le cube par automorphismes}

% archive sheet number: 22

